\scanpage{301}
\mspage{217}

4.9. \uline{Le théorème de \unsure{représentations indicielles|représentation indicielle}} \uline{pour} $\underline{\operatorname{Ind}}_{\pi}(C)$.

4.9.1. \unsure{Soient} $C$ une catégorie équivalente à une petite catégorie, $\pi_0$ un cardinal infini,

\[
M=\operatorname{Ind}_{\pi_0}(C)\subset C^\wedge,\tag*{(4.9.1)}
\]

$J$ une catégorie équivalente à une petite catégorie. Considérons l'inclusion \inserted{pl. fid} $C\hookrightarrow M$, et l'inclusion correspondante de $U$-catégories

\[
\underline{\operatorname{Hom}}(J,C)\hookrightarrow\underline{\operatorname{Hom}}(J,M).\quad\text{(pl. fid.)}.\tag*{(4.9.2)}
\]

On sait, \deleted{que} comme $C\subset M_{\pi_0}$, que les objets de $\underline{\operatorname{Hom}}(J,C)$ donnent des objets $\pi_0$-\uline{accessibles} de \deleted{$\unsure{C}$,} $\underline{\operatorname{Hom}}(J,M)$, du moins si on suppose

\[
\operatorname{card}\operatorname{red}J\leq\pi_0,\tag*{(4.9.\unsure{3|B})}
\]

ce que nous ferons désormais. (Cf 4.5.8, \unsure{1°|?°}, p.~119) Comme le deuxième \unsure{mb} de (4.9.2) satisfait $L_{\pi_0}$, ($M$ y satisfaisant), il s'ensuit que \deleted{l'in} le foncteur (4.9.2) se prolonge en

\[
\operatorname{Ind}_{\pi_0}(\underline{\operatorname{Hom}}(J,C))\hookrightarrow\underline{\operatorname{Hom}}(J,\underset{\substack{\Vert\\ \operatorname{Ind}_{\pi_0}(C)}}{M}),\tag*{(4.9.3)}
\]

et ce foncteur est \uline{pleinement fidèle} grâce au résultat d'accessibilité qu'on vient d'invoquer. Je m'intéresse à établir des cas où ce foncteur est essentiellement surjectif, i.e. est une \uline{équivalence} de catégories. Notons que pour que \inserted{$\varphi:J\unsure{\longrightarrow|\text{---}}M$} appartienne à l'image essentielle, il f. et s. qu'il existe un foncteur

\[
\Phi:I\longrightarrow\underline{\operatorname{Hom}}(J,C)\qquad\unsure{\Biggl\lbrace|\Biggl\vert}\begin{aligned}&I\ \text{grd devant }\pi_0\\ &\text{(et OPS }I\ \\&\text{ordonné)}\end{aligned},\tag*{(4.9.4)}
\]
% archive sheet number: 300
